\section{Trace Analysis and Partition Selection}
\label{sec:guided-selection}

Let \(D=(S^+,S^-)\) be a finite labeled sample over a declared alphabet
\(\Sigma\). We seek one partition \(P\) of \(\Sigma\) with maximum
cardinality among the compatible members of a decidable family
\(\mathfrak P\). Equivalently, we minimize
\(c(P)=|\Sigma|-|P|\).
For a block \(B\), the projection \(\pi_B\) erases letters outside \(B\)
and preserves the order of the remaining letters.
In the absence of state bounds, a fixed partition is compatible exactly
when for every \(y_i\in S^-\) some block \(B\in P\) satisfies
\(\pi_B(y_i)\notin\pi_B(S^+)\). Necessity follows because every component
must accept its positive images. Sufficiency follows by taking the finite
local languages \(\pi_B(S^+)\). We use this criterion for candidate
testing and the refinement property of projection equality for
candidate generation. With state bounds, the criterion remains necessary.

\subsection{Trace analysis and the feedback interface}

We retain the declared symbol order, the original trace order and a
stable index for each trace. In particular, unused symbols remain in the
alphabet and repeated occurrences of a labeled trace retain their
indices. A word occurring with both labels certifies infeasibility.
For \(\Sigma=\varnothing\), the empty conjunction accepts its sole
possible trace \(\epsilon\), so any negative trace certifies
infeasibility. In the absence of these two cases, analysis constructs
a starting partition and provides access to positive witnesses.

For each negative \(y=uabv\) with \(a\ne b\), analysis checks whether
the adjacent swap \(ubav\) is in \(S^+\). Each observed opposite-label
swap can supply an edge \(\{a,b\}\) with the two trace indices and the
swap position as evidence. To form \(P_0\), we keep the connected
components of the witness graph, including isolated symbols. The
implementation retains only edges that join two previously distinct
components. With preprocessing disabled, \(P_0\) is the singleton
partition.

\begin{lemma}[Mandatory components]
Every compatible partition, with or without state bounds, coarsens \(P_0\).
\end{lemma}
\begin{proof}
If \(a\) and \(b\) belonged to different blocks, projection onto any
block would erase at least one of them. Then \(uabv\) and \(ubav\)
would have identical projections on every block. Every product of
local languages would assign them the same label, in contradiction to
the input. Thus \(a,b\) share a block. The same-block relation is
transitive, so every component lies within one block.
\end{proof}

\paragraph{Procedure 1 (Analyze traces).}
\begin{enumerate}
\item Index the labeled traces over \(\Sigma\). Report a conflicting-label
  witness or an empty-alphabet negative when applicable.
\item Scan negative traces for adjacent swaps in the positive index.
  Join the two symbols' components for each witnessed edge and record
  the edge's evidence whenever a join occurs.
\item Return \(P_0\), the swap evidence, and access to positive matches
  on requested sets of symbols.
\end{enumerate}

For any proposed block \(B\) and negative index \(i\), a match query
returns a positive index \(j\) satisfying
\(\pi_B(x_j)=\pi_B(y_i)\), or reports that no such positive exists.
The implementation computes these queries from the ordered input words
and caches a row of negative-index matches per block. A projection failure
returns one negative index that has a match on every block of the current
partition. Selection can also query matches on unions of current blocks.

The last capability matters because projection feedback on current
blocks can lose cross-block order. With positive \(abc\) and the singleton
partition, compare the one-negative samples with \(bac\) and \(acb\),
respectively. In both samples every singleton's negative image equals its
positive image. The feedback tables contain the same fixed zero column
on every block, with the same positive witness. Yet the only compatible
single merge is \(\{a,b\}\) in the first sample and \(\{b,c\}\) in the
second. The example compares the singleton feedback directly, before
using any swap analysis.

\subsection{Candidate generation from a failed partition}

Suppose \(P\) fails the projection test on \(y_i\). Starting at \(P\),
merge two temporary blocks whenever the union still has a positive
match for that same \(y_i\). Stop after the first successful merge,
or continue until no pair can be merged with this property.
The resulting partition \(R\) coarsens \(P\), and every block of \(R\)
has its own positive witness for \(y_i\). The zero-step choice \(R=P\)
is also valid. These merges form the rejection evidence. Actual
children of \(P\) are produced by a single merge of two blocks of
\(P\) lying in distinct blocks of \(R\).

\begin{lemma}[Validity of rejection evidence]
Every refinement of \(R\) fails the projection test on \(y_i\).
\end{lemma}
\begin{proof}
For every block \(A\) of a refinement, choose the block \(C\in R\)
containing \(A\) and its recorded witness \(x_j\).
Since \(\pi_A\circ\pi_C=\pi_A\), equality of the images on \(C\)
implies equality on \(A\). Thus \(y_i\) has a positive match on every
block of the refinement.
\end{proof}

\begin{lemma}[Coverage by the generated children]
If \(Q\) is compatible and coarsens a projection-failing \(P\), then
one of the generated children of \(P\) refines \(Q\).
\end{lemma}
\begin{proof}
The preceding lemma implies that \(Q\) cannot refine \(R\).
A block of \(Q\) therefore intersects at least two blocks of \(R\).
Since both \(Q\) and \(R\) coarsen \(P\), that block of \(Q\) contains
two whole blocks of \(P\) lying in distinct blocks of \(R\). Merging
them gives a generated child still refining \(Q\).
\end{proof}

\paragraph{Procedure 2 (Select a partition).}
\begin{enumerate}
\item Initialize the frontier to \(\{P_0\}\).
\item For each partition in the current frontier, check membership in
  \(\mathfrak P\). Mark nonmembers for expansion by all pairwise merges
  and skip their feasibility tests.
\item For each member, run the projection test. Mark a projection failure
  with its negative index for witnessed expansion. If the test passes
  and there is no bound, return the partition. If bounds apply, use the
  bounded test on this partition. Return on its success and mark a
  definitive failure for expansion by all pairwise merges. Propagate
  an incomplete test as an incomplete search.
\item If the level has no success, expand its marked partitions using
  their prescribed rules. Deduplicate the next frontier. Repeat from
  step 2 if it is nonempty. Otherwise report infeasibility.
\end{enumerate}

The bounded test is an external decision procedure for the fixed
partition and its supplied bounds. A definitive bounded failure licenses
skipping that candidate. All pairwise merges are retained in this case.
Similarly, an inadmissible intermediate partition can lead to an
admissible candidate and provides no reason to restrict its children.

\begin{theorem}[Correctness and optimality]
For a finite alphabet and a decidable admissible family, the completed
unbounded procedure returns a compatible member of maximum cardinality,
or certifies that no compatible member exists. With a terminating,
sound and complete bounded test, the same statement holds for
compatibility under the supplied bounds.
\end{theorem}
\begin{proof}
Each level has finite size. One merge decreases the number of blocks
by one, so there are finitely many levels. For a compatible admissible
target \(Q\), the mandatory-components lemma ensures that \(P_0\)
refines \(Q\). If a strict refinement of \(Q\) in a frontier fails
the projection test, the coverage lemma supplies a child refining
\(Q\). If it is outside the family or fails the bounded test, merging
two of its blocks within one block of \(Q\) again supplies such a child.
Deduplication preserves that child as a member of the frontier.
Consequently, unless selection has already succeeded, \(Q\) is reached
at its own level. A returned partition is a member and passes the exact
test for the chosen regime. A compatible member with more blocks would
have been reached and accepted at an earlier level. Exhaustion with no
response would contradict reachability of any compatible member.
The initial direct infeasibility cases follow from identical words
requiring different labels, or from the empty conjunction.
\end{proof}

\subsection{Independent replay and search cost}

The unbounded certificate consists of the swap witnesses and records
for the completed levels. A projection record includes \(R\), the
negative index and one positive witness per block. A nonmember record
names the skipped partition. Replay checks the swaps from the input,
recomputes \(P_0\), validates records on the entire reconstructed frontier
and derives the next frontier independently of the search implementation.
It tests the final answer's membership and positive-image criterion.
For infeasibility by exhaustion, it requires a complete record of every
reached frontier and an empty next frontier. These checks establish the
hypotheses of the coverage proof directly. Under bounds, replay also
requires the external test to confirm each bounded rejection and the
answer's acceptance, so the guarantee is conditional on that test.

At a projection-failing partition with \(k\) blocks, let \(t_C\) count
the blocks of \(P\) contained in \(C\in R\). The number of generated
children before deduplication is
\[
  \binom{k}{2}-\sum_{C\in R}\binom{t_C}{2}.
\]
Each omitted pair is counted once within its containing block of \(R\).
This is a reduction in generated children. The implementation still
scans all pairs at each expanded partition. The greedy construction of
\(R\) makes \(O(k^3)\) match lookups in the worst case, since each stage
may scan quadratically many pairs before a merge and there are at most
\(k-1\) merges.

Let \(n=|\Sigma|\), let \(W\) be the number of words and let \(L\) be
their total length. On hash tables with expected constant-time operations
after hashing keys, the initial positive index costs \(O(W+L)\).
The simple swap routine can copy and hash a word at every candidate
position, giving \(O(\sum_{y\in S^-}|y|^2)\) time for that work.
Its representative walks cost \(O(n(L+n))\) as a conservative upper
bound, including the final component pass. A missing cached projection
row costs \(O(n+W+L)\) expected time and stores one match per negative.
Total cost further includes frontier sorting and deduplication,
candidate tests, per-pair work, certificate output and independent
replay. There can be exponentially many blocks and partitions to
consider. The child-count identity alone gives neither an end-to-end
speed guarantee nor a lower bound on other exact algorithms.

\subsection{Evaluation}

We use exhaustive partition enumeration as a correctness reference on
small inputs. All 20,232 results and certificate replays agreed with
the reference across 2,529 inputs and eight configurations. The test
set includes exhaustive labelings over alphabets of at most two
symbols, seeded samples over three to six symbols, and targeted cases.
Scripted bounded replies exercise the selector's interface.
The timing comparison holds the projection test and the
replay budget constant. It varies two choices independently, namely
mandatory-component preprocessing and witnessed expansion with greedy
pair merges to exhaustion. Table~\ref{tab:guided-evidence} reports the
number of completed, independently verified runs out of four and
their median search-to-verification time in milliseconds. The inputs
are synthetic. Differences in completion or median time apply to
the tested inputs and budget.

\begin{table}[tb]
\caption{Completed runs and median total time in milliseconds in parentheses.
Each run had a total budget of 0.5 seconds.}
\label{tab:guided-evidence}
\centering
\small
\setlength{\tabcolsep}{4pt}
\begin{tabular}{lrrrr}
\hline
& \multicolumn{2}{c}{Preprocessing off} & \multicolumn{2}{c}{Preprocessing on} \\
Sample & All pairs & Guided & All pairs & Guided \\
\hline
Manual example & 4/4 (0.30) & 4/4 (0.28) & 4/4 (0.27) & 4/4 (0.32) \\
Singleton optimum (12) & 4/4 (0.22) & 4/4 (0.20) & 4/4 (0.21) & 4/4 (0.20) \\
Greedy example (4) & 4/4 (0.31) & 4/4 (0.35) & 4/4 (0.20) & 4/4 (0.19) \\
One group (6) & 4/4 (31.17) & 4/4 (11.40) & 4/4 (34.30) & 4/4 (12.17) \\
Two groups (6) & 4/4 (29.54) & 4/4 (8.69) & 4/4 (39.00) & 4/4 (8.79) \\
Two groups (8) & 0/4 & 4/4 (111.85) & 0/4 & 4/4 (94.41) \\
Three groups (9) & 0/4 & 4/4 (151.84) & 0/4 & 4/4 (136.64) \\
Three groups (12) & 0/4 & 0/4 & 0/4 & 0/4 \\
\hline
\end{tabular}
\end{table}



In the six-symbol two-group case with preprocessing off, guided
expansion reduced the number of tested partitions from 174 to 42.
In the four-symbol case, preprocessing alone reduced that count from
two to one, while guided expansion alone had a slightly higher median
time. Both guided variants completed the eight- and nine-symbol
grouped cases within the budget. All four variants remained incomplete
on the twelve-symbol three-group case.
